\section{Limits of Orderbook-Only Inference}\label{sec:limits}

\subsection{Sign-agreement against on-chain ground truth}\label{ssec:sign-agreement}

Six trade-based measures (effective spread, realized spread, Roll,
Abdi-Ranaldo, Kyle's $\lambda$, Amihud) require an aggressor sign for
each trade. Standard practice in equity microstructure infers that
sign from a quote-driven feed via Lee-Ready, which assumes the feed
exposes enough information to distinguish buyer-initiated from
seller-initiated trades.

Polymarket's public WebSocket feed does not. We infer trades from
the feed under a LOOSE rule (every resting-size decrement counts) and
match the inferred buckets against on-chain \texttt{OrderFilled}
events at 5\,s and exact-price granularity. We run this matching over
four disjoint 7-day windows (2026-02-28~--~03-06, 03-07~--~03-13,
03-14~--~03-20, 03-21~--~03-27) on the top-100 stratum, computing
sign-agreement and a 95\% bootstrap CI for cells with at least ten
matched buckets.

\begin{itemize}
  \item Panel mean (109 valid cells of 400): \textbf{0.615}
  \item Volume-weighted by matched-bucket count (total 82{,}544):
        \textbf{0.586}
  \item Panel median 0.591; IQR $[0.526, 0.681]$;
        $p_{10}$--$p_{90}$ $[0.464, 0.841]$
  \item Per-window medians 0.586 / 0.591 / 0.618 / 0.645
\end{itemize}

A volume-weighted sign-agreement of $\sim 59\%$ sits just above the
50\% chance baseline. Even when an inferred bucket matches an
on-chain bucket in time and price, the inferred aggressor direction
is wrong about two trades in five. The mechanism is the feed itself:
\texttt{price\_change} updates broadcast a post-match snapshot of the
resting book without identifying the taker. The \texttt{change\_side}
field marks which side of the book moved, not which side initiated
the trade, and using it as a sign proxy produces the $\sim 59\%$
agreement rate.

\subsection{Propagation to direction-dependent measures}\label{ssec:strict-vs-onchain}

A noisy sign propagates to every measure that consumes it. We compute
all six trade-based measures under STRICT inference and under
authoritative on-chain trades on the top-100 stratum over a 7-day
window (2026-03-07~--~03-14); the window is shorter than the full
scrape because STRICT inference is $O(n^2)$ in the recent-event list
and the 28-day version did not finish in tractable wall time on this
panel size.

Three patterns emerge from the panel distribution.

\textit{Trade-count divergence is direction-unpredictable.}
Among 31 markets where both inference paths yielded non-empty trade
streams, the STRICT/on-chain count ratio has a median of 0.23
($p_{10}$--$p_{90}$ $[0.01, 1.45]$; min--max $[0.004, 3.62]$).
STRICT typically under-counts by an order of magnitude; on a small
tail of markets it over-counts by up to $3.6\times$. There is no
correctable scaling factor.

\textit{Effective spread sign-flips on two-thirds of the comparable
panel.} Among 24 markets with both sources non-empty,
16 (67\%, Wilson 95\% CI $[0.47, 0.82]$) have STRICT and on-chain
effective half-spreads of opposite sign. The STRICT median is
$+0.0048$\,prob~pp; the on-chain median is $-0.000075$\,prob~pp.
Inferred trade directions induce a positive spread that the
authoritative trade record does not support.

\textit{Kyle's $\lambda$ sign-flips on three-fifths of the comparable
panel.} Of 25 markets with both sources non-null and non-NaN, 15
(60\%, Wilson 95\% CI $[0.41, 0.77]$) have STRICT and on-chain
Kyle's $\lambda$ of opposite sign.

These distributions replace the n=3 comparison reported in earlier
versions of this paper. The artifact
\texttt{artifacts/measures\_compare\_top100.parquet} contains the
per-market panel. We note that only 24--31 of 100 top markets had
both sources non-empty in the 7-day slice (the remainder had zero
events on one side); the comparable-market subset is what supports
the sign-flip rates above. Restricting the effective-spread
comparison to markets with at least 100 STRICT trades raises the
sign-flip rate from 67\% to 73\% (16/22, Wilson 95\% CI
$[0.52, 0.87]$); the Kyle's $\lambda$ rate is essentially unchanged
(59\% at the 100-bucket threshold). Small-sample selection on
n_{\mathrm{strict}}$ does not generate the result.

\subsection{Implications for Polymarket microstructure research}\label{ssec:methodology}

Polymarket microstructure results that depend on trade direction
require on-chain \texttt{OrderFilled} events as the direction source,
not the public WebSocket feed. We release a replication package
(\texttt{polydata.\allowbreak onchain.\allowbreak trades.\allowbreak load\_onchain\_trades})
that performs the off-chain to on-chain join, together with a small
patch set
(\texttt{polydata.\allowbreak measures.\allowbreak \_trade\_source.\allowbreak resolve\_trades})
that lets existing measure code accept injected authoritative trades
without a rewrite.

\subsection{Methodology caveats}\label{ssec:methodology-caveats}

Three places where our compute approach trades accuracy for
tractability on a 28-day, 600-market panel.

\textit{Sample step.}
Quote samples and bucketed lookups use a 60-second grid for the panel
compute. Sensitivity tests on the top-5 markets across
$\{1, 10, 60, 300\}$\,s (artifact
\texttt{sample\_step\_sensitivity.parquet}) show that Roll is
invariant by construction, that effective spread and Amihud are
stable to within $\sim 10$--20\% at 60\,s relative to 1\,s on most
markets, and that Kyle's $\lambda$ varies by orders of magnitude
across step sizes on noisy markets. Kyle's $\lambda$ is the most
fragile estimator at any sample step, consistent with its
sign-flip behaviour in Section~\ref{ssec:strict-vs-onchain}.

\textit{Cross-sectional SF8.}
SF8 (Section~\ref{ssec:sf8}) is cross-sectional, not within-market
over time. A within-market temporal regression would require
per-market depth time-series that the current panel does not
materialise.

\textit{Wash-filter lower bound.}
SF7 (Section~\ref{ssec:sf7}) returns a lower bound by construction;
the gap to the network-classifier estimates documented in
\citet{cong-li-niu-2023} for cryptocurrency venues is the share of
wash-like activity that requires multi-counterparty graph
classification rather than direct or one-step roundtrip detection.

\textit{MEV as a confound for the buyer/seller asymmetry.}
Polymarket settles on Polygon, whose mempool is publicly accessible
to bots that can sandwich incoming orders, place runner trades, or
execute related back-running strategies. We do not analyse mempool
data in this paper, so the on-chain trade record we treat as ground
truth includes whatever fills MEV bots produce on top of organic
flow. The 60/40 seller/buyer asymmetry we observed in early panel
work could reflect retail aggressor behaviour, MEV bot activity, or
some mixture of the two. A within-market decomposition would
require either Polygon mempool capture during the scrape window or
a wallet-clustering approach against known MEV searchers; both are
out of scope here and flagged for future work.
